\chapter{DishBrain}
\chaptermark{DishBrain}
\label{sec:dishbrain}

\section{Mạng trong đĩa nuôi chơi quần vợt}

DishBrain là tên các tác giả đặt cho bàn thử, trên đó một mẫu nuôi cấy nơ-ron sống, mọc trên con chip có điện cực, chơi một trò quần vợt đơn giản với một cây vợt~\cite{Kagan2022}. Đây là thí nghiệm được bàn luận nhiều nhất với cỗ máy từ nơ-ron sống (tổng quan về những cỗ máy như vậy nằm ở chương~\ref{sec:livingmachine}), và đặc biệt tiện cho cuốn sách này. Ở đây, cả mạng nhỏ đều nằm trong tầm tay: mọi đầu vào do người làm thí nghiệm đặt, mọi đầu ra đều được ghi. Mọi câu hỏi của các chương trước, như tín hiệu được mã hóa ra sao, ai dạy, đầu dò thấy gì, đều có thể đặt cho một mạng không có mắt, không có cơ và không có nơ-ron \nt{DOPA}. Ta sẽ mổ xẻ thí nghiệm như kỹ sư mổ xẻ bàn thử của người khác: sơ đồ, đầu vào, đầu ra, thầy dạy, kết quả, tranh cãi.

\section{Bàn thử}

Nơ-ron được lấy từ vỏ não phôi chuột, khoảng $800\,000$ tế bào mỗi chip, hoặc được nuôi từ tế bào gốc người, khoảng một triệu mỗi chip. Trong vài tuần, chúng lớn lên trên chip, mọc đan vào nhau và bắt đầu tự phát ra xung và chùm xung. Trên một diện tích $8$~mm$^2$ bên dưới chúng có $26\,000$ điện cực bạch kim; có thể ghi đồng thời tối đa $1024$ điện cực, còn kích thích trên thực tế thì qua tám điện cực.

Quanh mẫu nuôi cấy là một vòng kín (hình~\ref{fig:dishloop}). Máy tính mô phỏng trò chơi: bóng bay, dội vào tường, vợt di chuyển lên xuống. Vị trí bóng được chuyển thành dòng điện trên tám điện cực ``cảm giác''. Xung từ hai vùng ``vận động'' được đếm theo cửa sổ $10\ms$ và đẩy vợt. Sau mỗi cú đỡ trúng hoặc trượt, mẫu nuôi cấy nhận thêm một kích thích nữa, tức phản hồi. Cửa sổ $10\ms$ là nhịp của máy tính trong bàn thử, không phải của mẫu nuôi cấy: bản thân mạng vẫn làm việc bất đồng bộ, như mọi mạng trong cuốn sách này.

\begin{figure}[t]
\centering
\begin{tikzpicture}[>=Stealth,
  blk/.style={draw,very thick,rounded corners=3pt,align=center,font=\small,minimum height=1.2cm},
  lab/.style={font=\scriptsize,align=center}]
\node[blk,fill=blue!6,text width=3.4cm] (C) at (0,0) {mẫu nuôi cấy trên chip\\\scriptsize $\sim10^6$ nơ-ron, $26\,000$ điện cực};
\node[blk,fill=orange!10,text width=3.0cm] (G) at (7.2,0) {máy tính:\\trò quần vợt};
\draw[->,very thick,blue!60!black] (G.north west) to[bend right=25] node[lab,above] {bóng: \emph{vị trí} --- điện cực nào trong 8,\\\emph{tần số} --- 4--40 xung/s} (C.north east);
\draw[->,very thick,red!70!black] (C.south east) to[bend right=25] node[lab,below] {vợt: số xung của các vùng\\``lên'' và ``xuống'' trong $10\ms$} (G.south west);
\draw[->,thick,densely dashed,green!45!black] (G.east) -- ++(0.9,0) |- ++(0,-2.6) -| node[lab,pos=0.25,below] {trúng: kích thích đoán trước được; trượt: không đoán trước được} (C.south);
\end{tikzpicture}
\caption{Vòng lặp DishBrain. Mũi tên xanh lam là đầu vào: vị trí bóng được ghi bằng mã vị trí và mã tần số (chương~\ref{sec:code}). Mũi tên đỏ là đầu ra: hiệu số đếm của hai vùng đẩy vợt. Đường nét đứt xanh lục là phản hồi sau mỗi lượt đánh, ``thầy dạy'' duy nhất của bàn thử.}
\label{fig:dishloop}
\end{figure}

\section{Đầu vào và đầu ra}

\emph{Đầu vào} được ghi cùng lúc bằng hai mã của chương~\ref{sec:code}. Điện cực nào trong tám điện cực được kích thích cho biết bóng ở độ cao nào: đó là mã vị trí, như ở các cảm biến của chương~\ref{sec:sensors}. Xung kích thích đến dày hay thưa cho biết bóng còn xa bao nhiêu: $4$ xung mỗi giây khi bóng ở bức tường xa, và tăng tuyến tính tới $40$ xung mỗi giây khi bóng ở bức tường phía vợt. Đó là mã tần số; biên độ xung của bóng là $75$~mV. Mẫu nuôi cấy không ``biết'' trước mã nào: mã do người làm thí nghiệm đặt ra, và mạng phải học cách đọc nó.

\emph{Đầu ra} được đọc giống như ở các giao diện của chương~\ref{sec:debugging}. Xung của một vùng vận động đẩy vợt lên, của vùng kia đẩy xuống; ở dạng đơn giản nhất, với mỗi cửa sổ
\begin{empheq}[box=\keybox]{equation}
\Delta p \;=\; c\,\bigl(n_\uparrow - n_\downarrow\bigr),
\label{eq:dishreadout}
\end{empheq}
trong đó $n_\uparrow$, $n_\downarrow$ là số xung của hai vùng trong $10\ms$, còn $c$ là hệ số của bàn thử. Đây là mã hai dây của chương~\ref{sec:glutcomputer}: chiều chuyển động không nằm ở dấu của tín hiệu, mà ở chỗ dây nào trong hai dây ``to'' hơn.

Hãy xem nhiễu của đầu ra như vậy. Nếu một vùng phát tổng cộng $R$ xung mỗi giây, thì trong cửa sổ $T=10\ms$ trung bình có $RT$ xung, và theo~\eqref{eq:rateerror}, độ tản mạn tương đối bằng $1/\sqrt{RT}$. Với $R=1000$ xung mỗi giây, con số này khoảng $30$ phần trăm ở mỗi cửa sổ; với $R=100$ thì hơn một trăm. Vợt chắc chắn sẽ rung, và mạng chỉ có một cách để học: dịch hiệu số đếm \emph{trung bình} về đúng phía. Đây cũng là những định luật giới hạn mọi giao diện ở chương~\ref{sec:debugging}, chỉ khác là giờ chúng tác động ở cả hai phía của vòng lặp.

\section{Thầy dạy không có dopamine}

Điều thú vị nhất ở bàn thử là thầy dạy. Mẫu nuôi cấy nơ-ron vỏ não không có nơ-ron \nt{DOPA}, và bàn thử không gửi tín hiệu ``tốt hơn mong đợi'' nào cả. Phản hồi thì khác. Mạng đỡ trúng bóng thì cả tám điện cực cảm giác cùng nhận một kích thích ngắn \emph{đoán trước được}: xung $75$~mV, $100$ xung mỗi giây trong $0.1\,\text{s}$. Mạng đỡ trượt thì trong bốn giây có một kích thích mà các tác giả gọi là \emph{không đoán trước được}: xung $150$~mV, $5$ xung mỗi giây, vào những thời điểm ngẫu nhiên, trên các điện cực chọn ngẫu nhiên; sau đó là bốn giây nghỉ không kích thích, rồi bóng xuất phát lại~\cite{Kagan2022}.

Các tác giả xuất phát từ giả thuyết rằng mạng hành động sao cho đầu vào của nó trở nên dễ đoán~\cite{Friston2010}. Với kỹ sư, ý nghĩa rất đơn giản: mẫu nuôi cấy có đúng một cách để thoát khỏi bốn giây nhiễu ngẫu nhiên, đó là đỡ trúng bóng. Ta đã gặp ý tưởng này ở việc tiết kiệm xung trong chương~\ref{sec:code} và ở dự đoán khung hình trong chương~\ref{sec:recognition}.

Nhưng có thật cần đến tính dễ đoán ở đây không? Cũng có thể có một cơ chế thô hơn. Giả sử kích thích không đoán trước được sau cú trượt chỉ đơn giản \emph{xáo trộn} những thay đổi gần đây trong mạng, còn kích thích êm sau cú trúng thì không đụng đến chúng. Hãy mô tả cấu hình của mạng bằng một vectơ $\boldsymbol\psi$, và giả sử ở cấu hình $\boldsymbol\psi$ mạng đỡ trượt với xác suất $P_{\text{trượt}}(\boldsymbol\psi)$. Nếu các cú hích là đối xứng, tức từ $\boldsymbol\psi$ sang $\boldsymbol\psi'$ có xác suất bằng chiều ngược lại, thì ở trạng thái dừng, dòng rời khỏi mỗi cấu hình cân bằng với dòng đi vào, và tỷ lệ thời gian mạng ở trong cấu hình $\boldsymbol\psi$ bằng
\begin{empheq}[box=\keybox]{equation}
\rho(\boldsymbol\psi) \;\propto\; \frac{1}{P_{\text{trượt}}(\boldsymbol\psi)} .
\label{eq:dishdensity}
\end{empheq}
Cấu hình trượt ít hơn mười lần thì được giữ lâu hơn mười lần. Không ai báo cho mạng biết phải đổi trọng số theo hướng nào: các cấu hình tốt chỉ đơn giản là ít bị phá hơn.

Chúng tôi đã kiểm tra điều này trên một mô hình, rút gọn mẫu nuôi cấy về hai con số: vợt đến điểm $a\,y+b$ khi bóng đến độ cao $y$, và đỡ trúng nếu lệch ít hơn $0.1$. Sau cú trượt, cả hai số nhận một cú hích ngẫu nhiên, sau cú trúng thì giữ nguyên. Qua $2000$ lượt đánh, tỷ lệ bóng đỡ trúng tăng từ $0.21$ ở hai trăm lượt đầu lên $0.38$ ở năm trăm lượt cuối (bốn mươi ``mẫu nuôi cấy'', hình~\ref{fig:dishmodel}). Nếu cú hích đến sau \emph{mỗi} lượt đánh, phản hồi không mang thông tin gì, và tỷ lệ ở quanh $0.12$; nếu hoàn toàn không có cú hích, nó đứng ở $0.19$. Thí nghiệm trên bàn thử phù hợp với điều này: mức tăng có ý nghĩa trong một phiên chỉ có khi phản hồi đầy đủ. Khi thay kích thích bằng im lặng, cuối phiên mẫu nuôi cấy vẫn chơi tốt hơn so với không có phản hồi, nhưng hiệu ứng nhỏ hơn~\cite{Kagan2022}; lưu ý rằng ngay trong điều kiện này, sau cú trượt là một khoảng nghỉ dài, còn sau cú trúng là khoảng nghỉ ngắn, nên sự khác biệt giữa hai kết cục không mất hẳn.

\begin{figure}[t]
\centering
\begin{tikzpicture}[x=1cm,y=5cm]
\draw[->,gray!70!black] (0,0) -- (10.6,0);
\draw[->,gray!70!black] (0,0) -- (0,0.5) node[above,font=\small,black] {tỷ lệ đỡ trúng};
\foreach \v in {0.1,0.2,0.3,0.4}{\draw[gray!60!black] (-0.06,\v) -- (0.06,\v); \node[left,font=\scriptsize] at (0,\v) {\v};}
\foreach \x/\f/\l/\lab in {1.8/0.21/0.38/{nhiễu sau\\cú trượt},5.2/0.13/0.11/{nhiễu sau\\mỗi lượt đánh},8.6/0.19/0.19/{không nhiễu}}{
  \fill[blue!35] (\x-0.8,0) rectangle (\x-0.05,\f);
  \fill[red!60!black] (\x+0.05,0) rectangle (\x+0.8,\l);
  \node[below,font=\scriptsize,align=center] at (\x,-0.01) {\lab};
  \node[above,font=\scriptsize] at (\x-0.425,\f) {\f};
  \node[above,font=\scriptsize] at (\x+0.425,\l) {\l};}
\fill[blue!35] (7.4,0.47) rectangle (7.7,0.495); \node[right,font=\scriptsize] at (7.75,0.482) {200 lượt đầu};
\fill[red!60!black] (7.4,0.41) rectangle (7.7,0.435); \node[right,font=\scriptsize] at (7.75,0.422) {500 lượt cuối};
\end{tikzpicture}
\caption{Mô hình ``xáo trộn khi trượt'' (\texttt{dishbrain\_model.py}): tỷ lệ bóng đỡ trúng ở đầu và cuối $2000$ lượt đánh, trung bình trên bốn mươi mô hình mẫu nuôi cấy. Việc học chỉ diễn ra nếu cú hích đi sau cú trượt chứ không đi sau cú trúng, cũng như trong thí nghiệm, nơi mức tăng có ý nghĩa trong một phiên chỉ có khi phản hồi đầy đủ.}
\label{fig:dishmodel}
\end{figure}

\begin{keyblock}
\begin{proposition}[Thầy dạy bằng cách xáo trộn]
\label{prop:dishscramble}
Nếu thất bại làm mạng bị lắc, còn thành công để yên cho nó, thì mạng tự trôi về những cấu hình ít sai hơn: thời gian ở trong một cấu hình tỷ lệ nghịch với xác suất thất bại~\eqref{eq:dishdensity}. Việc học như vậy không cần tín hiệu thưởng, không cần dấu của nó, cũng không cần tính dễ đoán tự thân; chỉ cần phản hồi sau thành công và sau thất bại khác nhau. Sự thay đổi hành vi của mẫu nuôi cấy đã được đo; còn cơ chế nào hoạt động trong đĩa nuôi, dự đoán đầu vào hay xáo trộn, thì chưa được xác lập.
\end{proposition}
\end{keyblock}

Hãy so sánh với chương~\ref{sec:dofa}. Ở đó nhiễu cũng dùng để tìm kiếm, nhưng dopamine có dấu: sai số $\delta$ cho biết phải dịch trọng số về phía nào, và bước đi theo gradient~\eqref{eq:perturbation}. Ở đây không có dấu, chỉ có ``giữ lại'' hoặc ``lắc''. Cách tìm kiếm này thô hơn và chậm hơn, nhưng đòi hỏi ở mạng ít hơn: không cần nơ-ron \nt{DOPA}, không cần bộ phê bình, không cần vết kéo dài hàng giây.

\section{Đã đo được gì và tranh cãi về điều gì}

Mỗi ván chơi kéo dài $20$ phút; người ta so sánh $5$ phút đầu với $15$ phút cuối. Ở các mẫu nuôi cấy có phản hồi đầy đủ, chuỗi cú đỡ trúng dài ra, còn số lần trượt ngay lập tức giảm đi; ở các đối chứng không có tế bào, không có trò chơi và với ``vợt'' do máy tính điều khiển thì không thấy gì tương tự. Mẫu nuôi cấy từ tế bào người trung bình đỡ lâu hơn mẫu của chuột. Sự cải thiện đáng tin cậy về thống kê, trên chín mẫu chuột và mười một mẫu người, với hơn một trăm ván cho mỗi nhóm, nhưng không lớn: mạng không trở thành người chơi giỏi, nó chỉ tốt hơn ngẫu nhiên một chút. Sự cải thiện đáng tin cậy trong một phiên; việc học bền vững từ phiên này sang phiên khác qua vài ngày thì các tác giả không đạt được~\cite{Kagan2022}. Trong công trình tiếp theo của cùng nhóm, người ta thấy khi chơi, hoạt động của mẫu nuôi cấy tiến gần tới chế độ \emph{tới hạn}, tức ranh giới giữa tắt dần và trận thác, chính là kiểu sống ở bờ vực trong chương~\ref{sec:gaba}, và càng gần thì chơi càng tốt~\cite{Habibollahi2023}.

Tranh cãi chính không nằm ở con số, mà ở lời lẽ. Tiêu đề bài báo khẳng định rằng nơ-ron trong đĩa nuôi ``thể hiện tri giác'' (sentience), và một nhóm nhà nghiên cứu đã phản bác: thay đổi hành vi trong vòng kín chưa phải là trí tuệ hay cảm giác, và kết luận cần được phát biểu khiêm tốn hơn~\cite{Balci2023}. Từ góc nhìn kỹ thuật, còn hai câu hỏi mở, và cả hai đều về cơ chế. Thứ nhất: mạng có học không, tức trọng số của nó có thay đổi lâu dài không, hay phản hồi chỉ dịch chuyển trạng thái hiện tại của nó? Thứ hai: có thật cần đến tính dễ đoán, hay chỉ cần bất kỳ sự khác biệt nào giữa đáp ứng với thành công và với thất bại, như trong mệnh đề~\ref{prop:dishscramble}? Một bàn thử mà mọi điện cực đều trong tầm tay cho phép trả lời cả hai, và có lẽ đó là giá trị chính của nó.

\section{Đặc tả bàn thử: cách dựng lại}
\label{sec:dishspec}

Dưới đây là mọi thứ cần để dựng lại một bàn thử như vậy: thiết bị, vòng lặp, mã hóa đầu vào, đọc đầu ra, phản hồi và giao thức thí nghiệm. Mỗi tham số đều ghi rõ nguồn. ``Bài báo'' nghĩa là giá trị lấy từ văn bản công trình~\cite{Kagan2022}. ``Lựa chọn'' nghĩa là bài báo không có, và để dựng lại thì phải tự đặt; chúng tôi đề xuất giá trị mặc định và nói cách kiểm tra nó.

\subsection*{Thiết bị và mẫu nuôi cấy}

\begin{center}
\footnotesize
\setlength{\tabcolsep}{4pt}
\renewcommand{\arraystretch}{1.15}
\begin{tabular}{>{\raggedright}p{4.0cm}>{\raggedright}p{6.9cm}>{\raggedright\arraybackslash}p{1.6cm}}
\toprule
Tham số & Giá trị & Nguồn \\
\midrule
chip & ma trận MaxOne: $26\,000$ điện cực bạch kim trên diện tích $8$~mm$^2$ & bài báo \\
ghi & tối đa $1024$ kênh đồng thời, $20\,000$ mẫu mỗi giây & bài báo \\
kích thích & về kỹ thuật tối đa $32$ điện cực, trong thí nghiệm là $8$ điện cực độc lập & bài báo \\
tế bào & vỏ não phôi chuột; nơ-ron người từ tế bào gốc (hai cách nuôi); tế bào HEK293T (không phải nơ-ron) làm đối chứng & bài báo \\
cấy & khoảng $10^6$ tế bào mỗi chip & bài báo \\
tuổi mẫu nuôi cấy khi làm thí nghiệm & chuột: từ $\sim10$ ngày; người: $14$--$24$ ngày hoặc $\sim80$ ngày tùy cách nuôi & bài báo \\
\bottomrule
\end{tabular}
\end{center}

\subsection*{Vòng lặp và thời gian}

Vòng lặp chạy theo cửa sổ $10\ms$ ($200$ mẫu): trong mỗi cửa sổ, xung trên các điện cực của vùng vận động được đếm, trò chơi được cập nhật và kích thích tiếp theo được phát ra. Độ trễ từ một xung trong mẫu nuôi cấy đến kích thích đáp lại là khoảng $5\ms$. Khi chơi, tín hiệu của mỗi điện cực đi qua bộ lọc thông cao Bessel bậc 2 với tần số cắt $100$~Hz; giá trị tuyệt đối của tín hiệu được làm trơn bằng bộ lọc thông thấp Bessel bậc 1 ở $1$~Hz, và ngưỡng xung tỷ lệ với giá trị tuyệt đối đã làm trơn này. Ngưỡng sáu độ lệch chuẩn của nền được bài báo nêu cho các phép đo hoạt động riêng lẻ, không phải cho lúc chơi. Để kích thích không bị tính là đáp ứng của mẫu nuôi cấy, mọi tín hiệu bị xóa khi có hơn $15$ xung lớn, trên $75$~mV, đến cùng lúc. Tất cả đều theo bài báo.

\subsection*{Mã hóa quả bóng}

\begin{center}
\footnotesize
\setlength{\tabcolsep}{4pt}
\renewcommand{\arraystretch}{1.15}
\begin{tabular}{>{\raggedright}p{3.4cm}>{\raggedright}p{7.5cm}>{\raggedright\arraybackslash}p{1.6cm}}
\toprule
Tham số & Giá trị & Nguồn \\
\midrule
vùng cảm giác & $8$ điện cực, bố trí sao cho các điện cực kề nhau ứng với các vị trí bóng kề nhau & bài báo \\
mã vị trí & số thứ tự điện cực $k=1,\dots,8$ cho biết vị trí bóng so với tâm vợt & bài báo \\
tọa độ nào & vị trí theo chiều dọc của bóng; để dựng lại: so với vợt, $k=\lceil 8(y-p+\tfrac12)\rceil$ khi $y-p\in[-\tfrac12,\tfrac12]$ & bài báo; công thức là lựa chọn \\
mã tần số & tần số kích thích từ $4$ đến $40$~xung/s mang khoảng cách tới vợt & bài báo \\
chiều của thang & $4$~xung/s khi bóng ở bức tường đối diện, tăng tuyến tính tới $40$~xung/s khi bóng ở bức tường phía vợt: $f=4+36\,(1-x/L)$, $x$ là khoảng cách tới tường phía vợt, $L$ là chiều rộng sân & bài báo \\
dạng xung & xung chữ nhật hai pha ngắn: pha dương trước, pha âm sau & bài báo \\
biên độ & $75$~mV; được chọn để không buộc các tế bào đang bị ức chế phải phát xung & bài báo \\
độ dài các pha & không nêu; lấy thiết lập chuẩn của hệ thống kích thích và giữ nguyên suốt thí nghiệm & lựa chọn \\
\bottomrule
\end{tabular}
\end{center}

\subsection*{Đọc vợt}

Bài báo có hai vùng vận động: xung ở vùng thứ nhất đẩy vợt lên, ở vùng thứ hai đẩy xuống; đóng góp của các vùng được lấy trọng số để tính đến mật độ tế bào và mức hoạt động khác nhau~\cite{Kagan2022}. Công thức chính xác không được đưa ra. Để dựng lại, ta lấy quy tắc~\eqref{eq:dishreadout}, $\Delta p=c\,(n_\uparrow-n_\downarrow)$ cho mỗi cửa sổ $10\ms$, với trọng số cân bằng các vùng theo hoạt động trung bình lúc nghỉ (lựa chọn): $n_\uparrow$ và $n_\downarrow$ được chia cho số đếm trung bình mỗi cửa sổ của vùng mình, đo trước khi chơi. Hệ số $c$ được chọn sao cho chênh lệch một số đếm trung bình dịch vợt $1\%$ chiều cao sân mỗi cửa sổ (lựa chọn); khi đó, một vùng trội hơn liên tục sẽ đưa vợt qua cả sân trong $1$~s.

Vị trí các vùng trên chip không được cho bằng tọa độ trong văn bản bài báo; chỉ có sơ đồ. Để dựng lại, chỉ cần ba nhóm điện cực không chồng lên nhau: tám điện cực cảm giác ở giữa và mỗi bên một nhóm điện cực ghi, một nhóm ``lên'', một nhóm ``xuống'' (lựa chọn).

\subsection*{Trò chơi}

Kích thước sân, chiều dài vợt và tốc độ bóng không được nêu trong bài báo; chỉ nói rằng bóng bay với tốc độ không đổi và dội lại từ các bức tường. Để dựng lại (lựa chọn): sân $1\times1$, vợt ở tường bên phải dài $0.2$ (trúng khi $|y-p|<0.1$, như trong mô hình~\texttt{models/dishbrain\_model.py}), bóng băng qua sân trong $1$~s. Cú trượt mà trong lượt đánh không có lần đỡ trúng nào được tính là ``ace''; lượt đánh có hơn ba cú đỡ liên tiếp là lượt ``dài'' (bài báo).

\subsection*{Phản hồi}

\begin{center}
\footnotesize
\setlength{\tabcolsep}{4pt}
\renewcommand{\arraystretch}{1.15}
\begin{tabular}{>{\raggedright}p{2.6cm}>{\raggedright}p{8.3cm}>{\raggedright\arraybackslash}p{1.6cm}}
\toprule
Sự kiện & Đưa vào gì & Nguồn \\
\midrule
trúng & kích thích đoán trước được: cả $8$ điện cực cảm giác cùng lúc, $75$~mV, $100$~xung/s, $100\ms$; trong thời gian này, kích thích vị trí bóng thông thường bị ngắt & bài báo \\
trượt & kích thích không đoán trước được: $150$~mV, $5$~xung/s, $4$~s, vào các thời điểm ngẫu nhiên trên các điện cực chọn ngẫu nhiên trong số $8$; sau đó $4$~s không kích thích, và bóng xuất phát theo hướng ngẫu nhiên & bài báo \\
luật ngẫu nhiên & không nêu; để dựng lại: thời điểm xung là dòng Poisson với tần số trung bình $5$~xung/s & lựa chọn \\
\bottomrule
\end{tabular}
\end{center}

\subsection*{Giao thức và đối chứng}

Một phiên kéo dài $20$~phút; so sánh $5$ phút đầu với $15$~phút cuối (bài báo). Ba thước đo kết quả: độ dài trung bình của lượt đánh, tỷ lệ ace và tỷ lệ lượt dài. Các điều kiện thí nghiệm (bài báo):
\begin{itemize}
\item \emph{kích thích}: vòng lặp đầy đủ, như mô tả ở trên;
\item \emph{im lặng}: thay cho kích thích khi trúng hoặc trượt là cùng khoảng thời gian ấy không có kích thích nào, sau đó bóng xuất phát theo hướng ngẫu nhiên;
\item \emph{không phản hồi}: sau cú trượt, bóng chỉ dội lại và bay tiếp, kích thích vị trí bóng vẫn tiếp tục;
\item \emph{nghỉ}: trò chơi vẫn chạy và vợt di chuyển theo hoạt động, nhưng hoàn toàn không có kích thích;
\item \emph{đối chứng}: chip chỉ có môi trường nuôi; tế bào HEK293T; ``mẫu nuôi cấy trong máy tính'', tức mô phỏng thay cho tế bào.
\end{itemize}
Cỡ mẫu trong loạt chính: $9$ mẫu chuột ($101$ phiên), $11$ mẫu người ($138$), $6$ chip chỉ có môi trường ($80$), $20$ mẫu ở trạng thái nghỉ ($42$), $3$ mô phỏng ($38$), tổng cộng $399$ phiên. Trong loạt về phản hồi: $486$ phiên trong $3$ ngày, mỗi ngày $3$ phiên, với các điều kiện ``kích thích'', ``im lặng'' và ``không phản hồi'' ($n=27$, $17$ và $15$). Độ dài trung bình của lượt đánh tăng từ $5$ phút đầu đến $15$ phút cuối chỉ ở các mẫu sống. Trong loạt về phản hồi, mức tăng có ý nghĩa theo thời gian chỉ có ở điều kiện ``kích thích''; cuối phiên, điều kiện ``im lặng'' vẫn hơn ``không phản hồi'', nhưng với hiệu ứng nhỏ hơn, còn việc học bền vững từ phiên này sang phiên khác trong ba ngày thì không có~\cite{Kagan2022}.

\subsection*{Chu trình của bàn thử từng bước}

Cứ mỗi $10\ms$, máy tính của bàn thử làm như sau.
\begin{enumerate}
\item Đọc số xung $n_\uparrow$, $n_\downarrow$ ở hai vùng vận động trong cửa sổ vừa qua và chuẩn hóa chúng theo số đếm trung bình của vùng lúc nghỉ.
\item Dịch vợt một đoạn $\Delta p=c\,(n_\uparrow-n_\downarrow)$ và giới hạn nó trong mép sân.
\item Dịch bóng một bước; cho bóng dội lại từ tường trên, tường dưới và tường trái.
\item Nếu bóng đã tới tường phải: khi $|y-p|<0.1$ là trúng, số đếm của lượt đánh tăng một, phát kích thích trúng ($100\ms$); nếu không thì là trượt, ghi lại lượt đánh, phát kích thích trượt ($4$~s), sau đó nghỉ $4$~s, và bóng xuất phát lại.
\item Nếu không, chọn điện cực $k$ theo vị trí dọc của bóng và tần số $f$ theo khoảng cách tới tường phía vợt, rồi đưa vào điện cực $k$ các xung hai pha $75$~mV với tần số $f$ cho đến cửa sổ tiếp theo.
\end{enumerate}
Như vậy là đủ để dựng một bàn thử tương thích với bàn thử đã công bố và kiểm tra câu hỏi chính của chương: dạy mẫu nuôi cấy là tính dễ đoán, hay chỉ là sự khác biệt giữa đáp ứng với cú trúng và với cú trượt. Muốn vậy, nên thêm vào ba điều kiện của bài báo một điều kiện thứ tư mà bài báo không có: sau cú trượt là kích thích \emph{đoán trước được}, cùng độ dài và năng lượng với kích thích không đoán trước được. Nếu mẫu nuôi cấy vẫn học được, thì vấn đề nằm ở sự khác biệt, không phải ở tính dễ đoán.


\section{Tổng kết}

\begin{table}[t]
\centering
\footnotesize
\setlength{\tabcolsep}{4pt}
\renewcommand{\arraystretch}{1.15}
\begin{tabular}{>{\raggedright}p{2.1cm}>{\raggedright}p{4.2cm}>{\raggedright}p{1.9cm}>{\raggedright\arraybackslash}p{4.6cm}}
\toprule
Khối của bàn thử & Cấu tạo & Chương & Điều đã biết \\
\midrule
đầu vào & vị trí (điện cực nào trong 8) và tần số (4--40 xung/s) & \ref{sec:code}, \ref{sec:sensors} & do người làm thí nghiệm đặt \\
đầu ra & hiệu số đếm của hai vùng trong $10\ms$~\eqref{eq:dishreadout} & \ref{sec:glutcomputer}, \ref{sec:debugging} & do người làm thí nghiệm đặt; nhiễu theo~\eqref{eq:rateerror} \\
thầy dạy & kích thích đoán trước được sau cú trúng, không đoán trước được sau cú trượt; không có dopamine & \ref{sec:dofa} & mức tăng có ý nghĩa trong một phiên chỉ khi phản hồi đầy đủ; với im lặng thì hiệu ứng nhỏ hơn; không bền vững từ phiên này sang phiên khác \\
cơ chế & dự đoán đầu vào hoặc xáo trộn khi trượt~\eqref{eq:dishdensity} & \ref{sec:dofa}, \ref{sec:recognition} & chưa xác lập; cả hai đều giải thích được kết quả \\
chế độ của mạng & tiến gần chế độ tới hạn khi chơi & \ref{sec:gaba} & mối liên hệ với chất lượng chơi đã đo \\
``tri giác'' & khẳng định của tác giả & --- & chưa được chứng minh; bị phản bác \\
\bottomrule
\end{tabular}
\caption{DishBrain qua con mắt kỹ sư.}
\label{tab:dishbrain}
\end{table}

DishBrain là cỗ máy từ nơ-ron sống ở dạng đơn giản nhất: đầu vào được đặt bằng mã, đầu ra được đọc bằng phép đếm, thầy dạy là sự khác biệt giữa đáp ứng với thành công và với thất bại (bảng~\ref{tab:dishbrain}). Một mạng không có mắt, cơ và dopamine đã học chơi được một chút, và chỉ riêng điều đó đã biến nó thành bàn thử cho các ý tưởng của cuốn sách. Dù cơ chế hóa ra là gì, dự đoán, xáo trộn hay một thứ thứ ba, câu trả lời sẽ được tìm theo cách chương~\ref{sec:debugging} khuyên: bằng đầu dò, tín hiệu thử và việc tắt từng phần trên một cỗ máy mà cuối cùng ta nhìn thấy được trọn vẹn.

\section{Bài tập}

\begin{exercise}[\lvlA]
\label{ex:22:1}
\emph{Cần dịch hiệu số đếm bao nhiêu.} Hai vùng cùng phát $R_\uparrow+R_\downarrow=2000$ xung mỗi giây, các dòng xung là Poisson. (a)~Độ tản mạn tương đối của số đếm một vùng trong $10\ms$ bằng bao nhiêu khi $R=1000$ và $R=100$? (b)~Trong $1$~s bóng bay, các độ dịch~\eqref{eq:dishreadout} cộng dồn. Với $R_\uparrow-R_\downarrow$ nhỏ nhất bằng bao nhiêu thì độ dịch trung bình lớn gấp đôi độ tản mạn?
\end{exercise}

\begin{exercise}[\lvlA]
\label{ex:22:2}
\emph{Cần bao nhiêu lượt đánh để thấy việc học.} Các lượt đánh độc lập, tỷ lệ đỡ trúng có phân phối nhị thức. Mỗi giai đoạn trong hai giai đoạn cần bao nhiêu lượt để mức tăng tỷ lệ đỡ trúng vượt hai độ lệch chuẩn của hiệu: (a)~từ $0.20$ lên $0.25$; (b)~từ $0.21$ lên $0.38$, như trong mô hình?
\end{exercise}

\begin{exercise}[\lvlB]
\label{ex:22:3}
\emph{Suy ra~\eqref{eq:dishdensity}.} Cấu hình $\boldsymbol\psi$ chạy trên một tập hữu hạn; khi trượt (xác suất $P(\boldsymbol\psi)>0$), mạng chuyển sang $\boldsymbol\psi'$ với xác suất đối xứng $K(\boldsymbol\psi,\boldsymbol\psi')$, khi trúng thì giữ nguyên. (a)~Chứng minh rằng $\rho\propto1/P$ là phân phối dừng. (b)~Tỷ lệ trượt bằng bao nhiêu khi có hai cấu hình với $P=0.9$ và $P=0.1$?
\end{exercise}

\begin{exercise}[\lvlB]
\label{ex:22:4}
\emph{Vùng hấp thụ của mô hình.} Vợt đến $p=\min\bigl(1,\max(0,ay+b)\bigr)$, bóng đến độ cao $y\sim U(0,1)$, trúng nếu $|p-y|<h=0.1$. (a)~Chứng minh rằng khi $|b|<h$ và $|a-1+b|<h$ mạng đỡ trúng mọi quả bóng, và tìm diện tích vùng này. (b)~Tìm bằng số tỷ lệ đỡ trúng trung bình khi $a\sim U(-1,1)$, $b\sim U(0,1)$.
\end{exercise}

\begin{exercise}[\lvlB]
\label{ex:22:5}
\emph{Mô phỏng: hàng rào quanh các cấu hình.} Trong một bản sao của \texttt{dishbrain\_model.py}, cho $200$ mẫu nuôi cấy mỗi mẫu $20\,000$ lượt đánh. Tỷ lệ đỡ trúng trong $500$ lượt cuối là bao nhiêu? Còn nếu sau mỗi cú hích, cắt cấu hình vào hình chữ nhật $a\in[-1,2]$, $b\in[-1,1]$? Giải thích sự khác biệt bằng bài tập~\ref{ex:22:3}.
\end{exercise}

\begin{exercise}[\lvlB]
\label{ex:22:6}
\emph{Thiết kế mã đầu vào.} Bóng được đỡ trúng nếu sai số về độ cao nhỏ hơn $h=0.1$; mạng đọc đầu vào trong $T=0.25$~s, với tần số đến $r$ phân biệt được khoảng $\sqrt{rT}$ mức, kích thích không dày quá $40$ xung mỗi giây. (a)~Mã vị trí trên tám điện cực có đủ không? (b)~Có thể truyền độ cao bằng tần số trên một điện cực không?
\end{exercise}

\begin{exercise}[\lvlC]
\label{ex:22:7}
\emph{Nghiên cứu: dự đoán hay xáo trộn?} Tổng quát hóa mô hình lên $d$ số điều chỉnh được ($p=\sum_{i<d}a_iy^i$). Số lượt đánh cần để đạt tỷ lệ đỡ trúng $0.5$ tăng theo $d$ như thế nào khi xáo trộn và khi dùng quy tắc có dấu của sai số~\eqref{eq:perturbation}? Thí nghiệm nào trên bàn thử sẽ phân biệt được hai giả thuyết?
\end{exercise}
